\documentclass[10pt,a4paper]{article}
\usepackage[utf8]{inputenc}
\usepackage[T1]{fontenc}
\usepackage[ngerman]{babel}
\usepackage[margin=1.7cm,top=1.5cm,bottom=1.7cm]{geometry}
\usepackage{titlesec}
\titleformat*{\section}{\normalfont\large\bfseries}
\titleformat*{\subsection}{\normalfont\normalsize\bfseries}
\titlespacing*{\section}{0pt}{8pt}{3pt}
\titlespacing*{\subsection}{0pt}{6pt}{2pt}
\usepackage{amsmath,amssymb}
\usepackage{listings}
\usepackage{xcolor}
\usepackage{array}
\usepackage{fancyhdr}
\usepackage{parskip}
\usepackage{needspace}

\newcommand{\mcbox}{\raisebox{-1pt}{\fbox{\rule{0pt}{7pt}\rule{7pt}{0pt}}}\hspace{6pt}}
\newcommand{\antwortlinien}[1]{%
  \par\vspace{2pt}%
  \newcount\zeilen \zeilen=#1
  \loop
    \noindent\rule{\linewidth}{0.4pt}\par\vspace{11pt}%
    \advance\zeilen by -1
  \ifnum\zeilen>0 \repeat
}
\lstset{language=Python,basicstyle=\ttfamily\footnotesize,keywordstyle=\bfseries,
  commentstyle=\itshape\color{gray},showstringspaces=false,frame=single,
  framerule=0.3pt,xleftmargin=8pt,framexleftmargin=6pt,aboveskip=6pt,belowskip=6pt}

\pagestyle{fancy}
\fancyhf{}
\fancyhead[L]{\small MDT5/2 -- Session 3: Evaluierung}
\fancyhead[R]{\small Übungsaufgaben zur Klausurvorbereitung}
\fancyfoot[C]{\small Seite \thepage}
\renewcommand{\headrulewidth}{0.4pt}

\begin{document}

\begin{center}
  {\Large\bfseries Übungsaufgaben Session 3}\\[4pt]
  {\large Kapitel 5: Metriken, ROC/AUC, Datenaufteilung, Kreuzvalidierung -- Praktikum 03}\\[4pt]
  \small Stil der Übungssammlung MDT5/2 -- generierte Aufgaben, keine Originale
\end{center}

\vspace{4pt}

\section*{Aufgabe 1 -- Multiple Choice mit Begründung}
\emph{Markieren Sie jeweils die einzige richtige von drei Aussagen und begründen Sie Ihre Entscheidung.}

\subsection*{1.1 \ Für die Datenaufteilung gilt:}
\mcbox Trainings-, Validierungs- und Testdaten müssen disjunkt sein; die Trennung beginnt
bereits bei der Vorverarbeitung.

\mcbox Validierungs- und Testdaten dürfen sich überlappen, solange die Trainingsdaten getrennt sind.

\mcbox Testdaten können nach einem Fehlschlag beliebig oft wiederverwendet werden.

Begründung:
\antwortlinien{2}

\subsection*{1.2 \ Für die Accuracy gilt:}
\mcbox Sie ist bei stark unbalancierten Daten die aussagekräftigste Metrik.

\mcbox Bei unbalancierten Daten kann sie hoch sein, obwohl die seltene Klasse überhaupt
nicht erkannt wird.

\mcbox Sie berücksichtigt nur die richtig erkannten Positiven.

Begründung:
\antwortlinien{2}

\subsection*{1.3 \ Für den Bias-Variance-Tradeoff gilt:}
\mcbox Overfitting zeigt sich durch einen großen Unterschied zwischen Trainings- und
Validierungsergebnissen.

\mcbox Underfitting entsteht durch zu hohe Modellvarianz.

\mcbox Bias und Varianz lassen sich immer gleichzeitig minimieren.

Begründung:
\antwortlinien{2}

\subsection*{1.4 \ Für die AUC (Area under the ROC-Curve) gilt:}
\mcbox Eine AUC von 0{,}5 bedeutet perfekte Trennung.

\mcbox Die AUC gibt an, wie wahrscheinlich eine zufällig gewählte Anomalie-Instanz
einen höheren Score erhält als eine zufällig gewählte normale Instanz.

\mcbox Die AUC hängt von einem fest gewählten Klassifikationsschwellwert ab.

Begründung:
\antwortlinien{2}

\section*{Aufgabe 2 -- Metriken rechnen}

Ein Novelty-Detection-Verfahren wird auf 1000 Testsamples angewendet:
950 normale Samples und 50 Anomalien. Positiv-Klasse ist \glqq Anomalie\grqq.
Das Verfahren liefert:
\begin{itemize}
  \item 40 Anomalien werden richtig als Anomalie erkannt
  \item 10 Anomalien werden fälschlich als normal erkannt
  \item 50 normale Samples werden fälschlich als Anomalie erkannt
  \item 900 normale Samples werden richtig als normal erkannt
\end{itemize}

a) Stellen Sie die Confusion Matrix auf (TP, FN, FP, TN).

\vspace{2.5cm}

b) Berechnen Sie Accuracy, Balanced Accuracy, Precision, Recall und $F_1$.
\antwortlinien{5}

c) Welche der Metriken beschreibt die Leistung des Verfahrens hier am ehrlichsten und warum?
\antwortlinien{3}

d) In der Anwendung ist jede übersehene Anomalie sehr teuer, Fehlalarme sind billig.
Sollte bei einem $F_\beta$-Maß eher $\beta=0{,}5$ oder $\beta=2{,}0$ gewählt werden? Begründung.
\antwortlinien{2}

\section*{Aufgabe 3 -- ROC-Kurve}

a) Welche Größen werden bei einer ROC-Kurve gegeneinander aufgetragen, und was wird dabei
variiert?
\antwortlinien{2}

b) Skizzieren Sie qualitativ drei ROC-Kurven: perfektes Verfahren, brauchbares Verfahren,
reines Raten.

\vspace{4.5cm}

c) Wie kann eine ROC-Kurve genutzt werden, um aus Anomalie-Scores einen Schwellwert
für die Klassifikation abzuleiten?
\antwortlinien{3}

\section*{Aufgabe 4 -- Kreuzvalidierung}

a) Beschreiben Sie den Ablauf einer k-Fold Cross Validation inklusive vorheriger Abspaltung
der Testdaten.
\antwortlinien{4}

b) Wann ist eine Kreuzvalidierung einer einmaligen Aufteilung in Training/Validierung/Test
vorzuziehen?
\antwortlinien{2}

c) Was bedeutet \glqq stratified\grqq{} Cross Validation und wann ist sie wichtig?
\antwortlinien{2}

d) Warum ist die Standard-Kreuzvalidierung von scikit-learn für Novelty Detection ungeeignet,
und wie sieht die Lösung aus der Vorlesung aus (Verteilung der Normaldaten und der Anomalien
auf die Teilmengen)?
\antwortlinien{4}

\section*{Aufgabe 5 -- Code beurteilen: Datenaufteilung}

\emph{Beurteilen Sie die drei voneinander unabhängigen Ausschnitte auf Fehlerfreiheit in Bezug
auf die Datenaufteilung und begründen Sie jeweils.}

\Needspace*{6cm}
\begin{lstlisting}
# Variante A
features, labels = load_data()
selector = SelectKBest(k=20)                 # Merkmalsauswahl
selected = selector.fit_transform(features, labels)

train_x, test_x, train_y, test_y = train_test_split(selected, labels,
                                                    test_size=0.2)
clf = SVC()
clf.fit(train_x, train_y)
print(np.mean(clf.predict(test_x) == test_y))
\end{lstlisting}
\antwortlinien{3}

\Needspace*{7cm}
\begin{lstlisting}
# Variante B
features, labels = load_data()
train_x, test_x, train_y, test_y = train_test_split(features, labels,
                                                    test_size=0.2)

pipe = Pipeline([('scaler', StandardScaler()), ('clf', SVC())])
grid = GridSearchCV(estimator=pipe,
                    param_grid={'clf__C': np.logspace(-2, 2, 10)},
                    cv=5, refit=False)
grid.fit(train_x, train_y)

pipe.set_params(**grid.best_params_)
pipe.fit(train_x, train_y)
print(np.mean(pipe.predict(test_x) == test_y))
\end{lstlisting}
\antwortlinien{3}

\Needspace*{6cm}
\begin{lstlisting}
# Variante C  (Handschriften mehrerer Probanden, spaeter sollen
#              Handschriften NEUER Personen erkannt werden)
features, labels, person_ids = load_handwriting()
train_x, test_x, train_y, test_y = train_test_split(features, labels,
                                                    test_size=0.2,
                                                    shuffle=True)
\end{lstlisting}
\antwortlinien{3}

\vspace{10pt}
\begin{center}
\footnotesize\itshape
Nach Bearbeitung: Antworten an Claude schicken (Foto genügt) -- Korrektur mit Begründungs-Check.
\end{center}

\end{document}
